\clearpage
\chapter*{2025 Enero}
\addcontentsline{toc}{chapter}{2025 Enero}

\begin{enunciado}
    \hypertarget{2025 A 1}{}
    \textbf{1. (2,5 puntos)} Sea \((X,T)\) un espacio topológico. Demostrar que un subconjunto \(M\) de \(X\) es denso en \(X\) si y sólo si para todo abierto no vacío \(U\) de \(X\) se tiene que \(U \cap M \neq \varnothing\).
\end{enunciado}

\begin{enunciado}
    \hypertarget{2025 A 2}{}
    \textbf{2. (3,5 puntos)} En el conjunto \(\mathbb{Z}\) de los números enteros se define la topología \(T\) mediante
    \[
        T = \{M \subset \mathbb{Z} \mid (0 \notin M) \text{ o } (\mathbb{Z} - M \text{ es finito})\}.
    \]
    \begin{enumerate}[label=\alph*)]
        \item Estudiar si \((\mathbb{Z},T)\) es conexo.
        \item Estudiar si \((\mathbb{Z},T)\) es compacto. 
    \end{enumerate}
\end{enunciado}

\begin{enunciado}
    \hypertarget{2025 A 3}{}
    \textbf{3. (4 puntos)}
    \begin{enumerate}[label=\alph*)]
        \item Sea \(f : X \to Y\) una aplicación suprayectiva y supongamos que \(X\) tiene la topología trivial \(T = \{X,\emptyset\}\). Probar que si \(S\) es una topología en \(Y\), la aplicación \(f : (X,T) \to (Y,S)\) es continua si y sólo si la topología \(S\) es la topología trivial de \(Y\).
        \item Demostrar que si \((X,T)\) es un espacio topológico con la topología \(T\) inducida por una aplicación \(f : X \to (Y,T')\), entonces se verifica que para todo espacio topológico \((M,S)\) y cualquier aplicación \(g : (M,S) \to (X,T)\), la aplicación \(g\) es continua si y sólo si es continua la aplicación \(f \circ g : (M,S) \to (Y,T')\).
    \end{enumerate}
\end{enunciado}